\documentclass[11pt]{beamer}

\usepackage[utf8]{inputenc}
\usepackage[T1]{fontenc}
\usepackage[default]{opensans}
\usepackage[english]{babel}
\usepackage{hyperref}
\usepackage{tikz}
\usepackage[dvipsnames]{xcolor}
\usepackage{graphicx}
\usetikzlibrary{calc,patterns}
\usepackage{bm}

\hypersetup{
  colorlinks=true,
  urlcolor=blue,
  linkcolor=blue,
  citecolor=blue
}

\usetheme{default}
\usecolortheme{dolphin}
\useinnertheme{circles}

\title{Belief-Aware Drone Navigation Under Partial Observability}
\subtitle{Project Proposal}
\author[M. Osipova]{Mariia Osipova\\[3pt]\textit{mosipova@udesa.edu.ar}}
\institute{\\[12pt]Under the supervision of Juan Felipe Perdomo}
\date[Apr 2026]{April 2026}

%----------------------------------------------------------------------------------------
% sources
%----------------------------------------------------------------------------------------
\newcommand{\srcSuttonBarto}{\href{http://incompleteideas.net/book/the-book.html}{Sutton \& Barto (2018)}}
\newcommand{\srcKochenderferADM}{\href{https://algorithmsbook.com}{Kochenderfer, Wheeler \& Wray (2022)}}
\newcommand{\srcKaelbling}{Kaelbling, Littman \& Cassandra (1998)}
\newcommand{\srcSchulman}{\href{https://arxiv.org/abs/1707.06347}{Schulman et al. (2017)}}
\newcommand{\srcTobin}{Tobin et al. (2017)}
\newcommand{\srcHwangbo}{Hwangbo et al. (2017)}
\newcommand{\SourceLine}[1]{\vfill{\tiny\textcolor{gray}{Sources: #1}}}

%----------------------------------------------------------------------------------------

\begin{document}

%----------------------------------------------------------------------------------------
% TITLE SLIDE
%----------------------------------------------------------------------------------------

\begin{frame}
    \titlepage
\end{frame}

%----------------------------------------------------------------------------------------
% TABLE OF CONTENTS SLIDE
%----------------------------------------------------------------------------------------

\begin{frame}{Presentation Overview}
  \tableofcontents
\end{frame}

%------------------------------------------------

\section{Motivation}

%------------------------------------------------

\begin{frame}{Why this problem?}

The question behind this project is simple:

\medskip
\begin{quote}
How should a drone act when it does not know exactly where it is?
\end{quote}

\medskip
What made this question compelling to me was the combination of:
\begin{itemize}
  \item reinforcement learning as a way to learn behavior from interaction,
  \item robotics as a concrete testbed for decision-making,
  \item uncertainty as the condition that makes the problem genuinely hard.
\end{itemize}

The setting is autonomous quadrotor navigation with noisy and incomplete sensing.

\SourceLine{\srcSuttonBarto; \srcKochenderferADM}
\end{frame}

%------------------------------------------------

\begin{frame}{From control to partial observability}

My path toward this topic was gradual:
\begin{itemize}
  \item first, classical control: for example, a PID controller for altitude stabilization;
  \item then, reinforcement learning: letting the agent learn rather than hand-coding behavior;
  \item finally, partial observability: realistic flight means noisy, delayed, and incomplete sensor data.
\end{itemize}

\medskip
This changes the nature of the task:
\begin{block}{Key shift}
The drone should not act only on raw observations. It should act on what it believes about its state, together with how certain that belief is.
\end{block}

\SourceLine{\srcSuttonBarto; \srcKaelbling}
\end{frame}

%------------------------------------------------

\section{Problem Formulation}

%------------------------------------------------

\begin{frame}{MDP vs. POMDP}

In a standard RL setup, the agent observes the true state and learns a policy
\[
\pi(a \mid s).
\]

For a drone with imperfect sensors, that assumption is false.

\medskip
A better model is a POMDP:
\begin{itemize}
  \item the true state is hidden,
  \item the agent receives noisy observations $o_t$,
  \item the relevant object is a belief over possible states.
\end{itemize}

\medskip
So the central issue is not only \emph{what action to take}, but also \emph{what to believe about the current state}.

\SourceLine{\srcKochenderferADM; \srcKaelbling}
\end{frame}

%------------------------------------------------

\begin{frame}{Research question}

\begin{block}{Central question}
Can a reinforcement-learning agent navigate a quadrotor more effectively under noisy and incomplete sensor data by explicitly maintaining a belief over its own state, compared to an agent that treats raw observations as ground truth?
\end{block}

\medskip
More specifically, I want to test whether explicit uncertainty representation improves:
\begin{itemize}
  \item robustness under noise,
  \item stability during training,
  \item sample efficiency,
  \item generalization to unseen conditions.
\end{itemize}
\end{frame}

%------------------------------------------------

\section{Method}

%------------------------------------------------

\begin{frame}{Experimental setup}

The project will use a simulated quadrotor navigation task built with:
\begin{itemize}
  \item \texttt{gym-pybullet-drones} for the environment,
  \item Stable-Baselines3 for PPO,
  \item NumPy / SciPy for belief statistics,
  \item Matplotlib for trajectory and training plots.
\end{itemize}

\medskip
Baseline task:
\begin{itemize}
  \item fly from the origin to a fixed waypoint.
\end{itemize}

Controllers to compare:
\begin{itemize}
  \item PID,
  \item PPO with raw observations,
  \item PPO with belief-augmented observations.
\end{itemize}

\SourceLine{\srcSchulman; \srcHwangbo}
\end{frame}

%------------------------------------------------

\begin{frame}{Phase 1: Baseline under ideal conditions}

Goal: establish whether the RL pipeline works before introducing uncertainty.

\medskip
I will compare:
\begin{itemize}
  \item a classical PID controller with direct state access,
  \item a PPO agent trained from scratch on the same waypoint task.
\end{itemize}

Main outputs of this phase:
\begin{itemize}
  \item reward curves,
  \item successful trajectories,
  \item a clean baseline for later degradation studies.
\end{itemize}

\SourceLine{\srcSchulman}
\end{frame}

%------------------------------------------------

\begin{frame}{Phase 2: Introduce uncertainty}

I will corrupt the observation pipeline using three noise mechanisms:
\begin{itemize}
  \item Gaussian noise on observed position,
  \item observation dropout or repeated stale measurements,
  \item external wind disturbances.
\end{itemize}

\medskip
This phase is designed to produce degradation curves and answer:
\begin{quote}
How quickly does performance collapse as sensing becomes less reliable?
\end{quote}

\medskip
It also tests whether the plain RL policy has learned genuine robustness or simply overfit to clean observations.

\SourceLine{\srcKochenderferADM; \srcSchulman}
\end{frame}

%------------------------------------------------

\begin{frame}{Phase 3: Belief-augmented observations}

This is the core contribution.

\medskip
From a sliding window of recent observations, compute:
\[
\hat{\mu}_t \qquad \text{and} \qquad \hat{\sigma}_t,
\]
where:
\begin{itemize}
  \item $\hat{\mu}_t$ is a filtered state estimate,
  \item $\hat{\sigma}_t$ is a measure of confidence.
\end{itemize}

The policy then receives an augmented input:
\[
[o_t, \hat{\mu}_t, \hat{\sigma}_t].
\]

Hypothesis: the agent will learn to behave more cautiously when uncertainty is high and more decisively when it is low.

\SourceLine{\srcKaelbling; \srcSchulman}
\end{frame}

%------------------------------------------------

\begin{frame}{Phase 4: Generalization and transfer}

I will compare three training regimes:
\begin{itemize}
  \item a specialist trained only under low noise,
  \item a specialist trained only under high noise,
  \item a generalist trained with domain randomization across noise levels.
\end{itemize}

\medskip
Evaluation will be done on held-out noise conditions.

\medskip
Key question:
\begin{block}{Transfer question}
Does the belief-augmented generalist transfer better to unseen conditions than policies specialized for a single regime?
\end{block}

\SourceLine{\srcTobin}
\end{frame}

%------------------------------------------------

\section{Expected Results}

%------------------------------------------------

\begin{frame}{Expected outcomes}

\begin{itemize}
  \item Under ideal conditions, PPO should roughly match PID after sufficient training.
  \item PID should be more sample-efficient because it requires tuning, not learning.
  \item Under moderate and high noise, plain PPO may degrade sharply if it overfits to clean observations.
  \item Belief-augmented PPO should outperform plain PPO when sensing becomes unreliable.
  \item Domain-randomized training should improve transfer to unseen noise regimes.
\end{itemize}

\medskip
The main expected result is that even a lightweight belief signal can improve RL behavior under partial observability.
\end{frame}

%------------------------------------------------

\begin{frame}{Planned outputs and evaluation}

Concrete outputs:
\begin{itemize}
  \item training curves for all agent variants,
  \item trajectory plots for PID, plain PPO, and belief-aware PPO,
  \item RMSE degradation curves under varying noise levels,
  \item a train/test generalization matrix across noise regimes,
  \item a written POMDP framing of the task.
\end{itemize}

\medskip
Evaluation criteria:
\begin{itemize}
  \item navigation accuracy,
  \item success rate,
  \item robustness to uncertainty,
  \item generalization outside the training regime.
\end{itemize}
\end{frame}

%------------------------------------------------

\section{Discussion}

%------------------------------------------------

\begin{frame}{Why this matters}

This project is a small but concrete test of a broader question in reinforcement learning:

\medskip
\begin{quote}
What representations need to be explicit for an agent to act well under uncertainty?
\end{quote}

\medskip
If belief-aware inputs help, that would support two claims:
\begin{itemize}
  \item partial observability is a practical bottleneck, not just a theoretical detail,
  \item explicit uncertainty information can change the behavior an RL agent learns.
\end{itemize}

It also connects control, robotics, RL, and POMDP-based reasoning in one concrete experimental setting.

\SourceLine{\srcSuttonBarto; \srcKochenderferADM; \srcKaelbling}
\end{frame}

%------------------------------------------------

\begin{frame}{Limitations and extensions}

Current limitations:
\begin{itemize}
  \item the belief representation is only a lightweight approximation,
  \item PPO is only one possible RL algorithm,
  \item the project remains in simulation.
\end{itemize}

\medskip
Possible extensions:
\begin{itemize}
  \item compare with Kalman-filter-based state estimation,
  \item test recurrent policies such as LSTMs,
  \item explore model-based RL or explicit POMDP solvers,
  \item study sim-to-real transfer later on.
\end{itemize}
\end{frame}

%------------------------------------------------

\section{References}

%------------------------------------------------

\begin{frame}{Selected references}
\footnotesize
\begin{itemize}
  \item Sutton, R. S., \& Barto, A. G. (2018). \textit{Reinforcement Learning: An Introduction}.
  \item Kochenderfer, M. J., Wheeler, T. A., \& Wray, K. H. (2022). \textit{Algorithms for Decision Making}.
  \item Kaelbling, L. P., Littman, M. L., \& Cassandra, A. R. (1998). Planning and acting in partially observable stochastic domains.
  \item Schulman, J. et al. (2017). Proximal Policy Optimization Algorithms.
  \item Hwangbo, J. et al. (2017). Control of a quadrotor with reinforcement learning.
  \item Tobin, J. et al. (2017). Domain randomization for transferring deep neural networks from simulation to real world.
\end{itemize}
\end{frame}

%------------------------------------------------

\begin{frame}
  \centering
  \Huge Thank you
\end{frame}

\end{document}
